\documentclass[11pt,a4paper]{article}
\usepackage[T1]{fontenc}
\usepackage{lmodern}
\usepackage[margin=27mm]{geometry}
\usepackage{amsmath,amssymb,amsthm,mathtools}
\usepackage{microtype}
\usepackage{booktabs,array}
\usepackage{enumitem}
\usepackage[hidelinks]{hyperref}
\setlength{\parskip}{5pt plus 1pt}
\setlength{\parindent}{0pt}
\setlength{\emergencystretch}{3em}
\setlist{itemsep=2pt,topsep=3pt}
\raggedbottom
\newtheorem{theorem}{Theorem}[section]
\newtheorem{lemma}[theorem]{Lemma}
\newtheorem{corollary}[theorem]{Corollary}
\newtheorem{proposition}[theorem]{Proposition}
\theoremstyle{definition}
\newtheorem{definition}[theorem]{Definition}
\newtheorem{example}[theorem]{Example}
\theoremstyle{remark}
\newtheorem{remark}[theorem]{Remark}
\newcommand{\Z}{\mathbb{Z}}
\newcommand{\N}{\mathbb{N}}
\newcommand{\Q}{\mathbb{Q}}
\newcommand{\F}{\mathbb{F}}
\newcommand{\ps}[1]{#1[[x]]}
\newcommand{\coeff}[2]{[x^{#1}]\,#2}
\newcommand{\hl}[1]{\textbf{#1}}
\newcommand{\para}[1]{\par\medskip\noindent\textbf{#1}\enspace}
\hypersetup{pdftitle={Remainders modulo 4 for the equation A(x² + 2x³) = x A(x): a proof of a conjecture in OEIS A389536},pdfauthor={Dev D. Goyal},pdfsubject={OEIS A389536}}
\title{Remainders modulo $4$ for the equation $A(x^2+2x^3)=x\,A(x)$\\[6pt]\large A proof of a conjecture in OEIS A389536}
\author{Dev D. Goyal}
\date{October 2026}
\begin{document}
\maketitle
\begin{abstract}
There is exactly one power series $A(x)=x+2x^2+2x^3+8x^4+10x^5+\cdots$ that satisfies the equation
$A(x^2+2x^3)=x\,A(x)$. Its coefficients $a(n)$ form the sequence A389536 of the On-Line Encyclopedia of Integer
Sequences (OEIS), contributed by Paul D.~Hanna. Hanna conjectured that, for $n>1$, the coefficient $a(n)$ leaves
remainder $2$ on division by $4$ when $n-1$ is a power of $2$, and is divisible by $4$ otherwise. We prove this
conjecture. The key observation is that, once every coefficient is replaced by its remainder modulo $4$, the series
$A(x)$ becomes $x+2x\,C(x)$, where $C(x)=x+x^2+2x^3+5x^4+\cdots$ is the generating function of the Catalan numbers.
The conjecture then follows from the classical fact that the coefficient of $x^n$ in $C(x)$ is odd exactly when $n$ is a power of
$2$. The result has also been checked by computer and proved formally in the Lean proof assistant.
\end{abstract}
\begin{center}\begin{minipage}{0.88\linewidth}\small
\textbf{How this paper was produced.} This paper reports the result of an AI research effort directed by
Dev~D.~Goyal. ChatGPT (OpenAI), using the GPT-6 model with Pro reasoning and working autonomously from prompts by
Dev~D.~Goyal, found the proof, wrote its Lean~4 formalization, and wrote a first draft of the paper.
Dev~D.~Goyal chose the problems and asked for proofs, formal verification and revisions, but did not supply proof
ideas. Claude Opus~5.5, used with Claude Code, then checked the work independently: it rebuilt the formal proofs from
source and confirmed that they rely only on Lean's three standard axioms, recomputed the sequence from its definition
compared it with the published terms and tested every claim numerically, and read every step of the proofs. It also wrote the present
version of this paper.
Dev~D.~Goyal has reviewed the proofs and is responsible for them; this review is not an independent expert validation.
The paper has not been peer reviewed.
The section ``About this paper'' at the end gives the details.
\end{minipage}\end{center}

\section{Introduction}

\subsection{The sequence}\label{sec:sequence}
The \emph{On-Line Encyclopedia of Integer Sequences} (OEIS)~\cite{oeisdb} is a large public database of integer
sequences. Many
entries also record \emph{conjectures}: patterns that the contributor observed in the data but did not prove.

This paper is about entry A389536~\cite{oeis}, contributed by Paul~D.~Hanna in October 2025. The entry reads, in part:
\begin{quote}\small\ttfamily\raggedright G.f. A(x) satisfies: A(x) = A(x\^{}2 + 2*x\^{}3) / x.\par\smallskip The radius of convergence r of the g.f. A(x) is r = 1/2, with A(-r) = 0.\par\smallskip Conjecture: for n > 1, a(n) == 2 (mod 4) when n = 2\^{}k + 1 (for some k >= 0) otherwise a(n) is divisible by 4.\par\smallskip\hfill\textrm{--- OEIS A389536, Paul D. Hanna, October 2025}\end{quote}

Here ``G.f.''\ stands for \emph{generating function}: the terms $a(1),a(2),a(3),\dots$ of the sequence are the
coefficients of the power series $A(x)=a(1)\,x+a(2)\,x^2+a(3)\,x^3+\cdots$. In the equation, $A(x^2+2x^3)$ is the series
obtained by substituting the polynomial $x^2+2x^3$ for $x$ in $A$ (Section~\ref{sec:background} explains this). Multiplying both
sides by $x$, we write the equation without division:
\begin{equation}\label{eq:def}
A\bigl(x^2+2x^3\bigr)=x\,A(x).
\end{equation}
Equations like this one, which relate $A(x)$ to $A$ evaluated at a higher power of $x$ (here, essentially $x^2$), are
called \emph{Mahler-type equations}, after the mathematician Kurt Mahler, who studied them. Equation~\eqref{eq:def} is
\emph{linear}: if $A$ is a solution, so is every multiple of $A$. So one coefficient must be fixed, and the entry fixes
$a(1)=1$. Then all other coefficients are determined (Theorem~\ref{thm:main}). The first terms, starting at $n=1$, are
1, 2, 2, 8, 10, 12, 32, 80, 202, 356, 540, 944, \dots

The second line of the quotation is a statement about $A(x)$ as a function of a real or complex number $x$. This paper
treats $A$ only as a formal power series (a list of coefficients), and does not address that statement.

\subsection{The question}
Hanna's conjecture says that, for $n>1$, the remainder of $a(n)$ on division by $4$ is $2$ when $n$ is one more than a
power of $2$ (that is, $n=2,3,5,9,17,\dots$), and $0$ otherwise. The first row of the table shows the computed remainders of $a(n)$; the
indices $n=2^k+1$ are printed in bold. The second row anticipates the main result: it shows the remainders modulo $4$ of
the coefficients of $x+2x\,C(x)$, where $C(x)$ is the Catalan series introduced below.
\begin{center}\small\setlength{\tabcolsep}{3.2pt}
\begin{tabular}{>{\raggedright\arraybackslash}p{8.8em}>{\centering\arraybackslash}p{1.15em}>{\centering\arraybackslash}p{1.15em}>{\centering\arraybackslash}p{1.15em}>{\centering\arraybackslash}p{1.15em}>{\centering\arraybackslash}p{1.15em}>{\centering\arraybackslash}p{1.15em}>{\centering\arraybackslash}p{1.15em}>{\centering\arraybackslash}p{1.15em}>{\centering\arraybackslash}p{1.15em}>{\centering\arraybackslash}p{1.15em}>{\centering\arraybackslash}p{1.15em}>{\centering\arraybackslash}p{1.15em}>{\centering\arraybackslash}p{1.15em}>{\centering\arraybackslash}p{1.15em}>{\centering\arraybackslash}p{1.15em}>{\centering\arraybackslash}p{1.15em}}
\toprule
$n$ & 1 & \hl{2} & \hl{3} & 4 & \hl{5} & 6 & 7 & 8 & \hl{9} & 10 & 11 & 12 & 13 & 14 & 15 & 16 \\ \midrule
$a(n)\bmod 4$ & 1 & \hl{2} & \hl{2} & 0 & \hl{2} & 0 & 0 & 0 & \hl{2} & 0 & 0 & 0 & 0 & 0 & 0 & 0 \\
$\coeff{n}{(x{+}2xC)}\bmod 4$ & 1 & \hl{2} & \hl{2} & 0 & \hl{2} & 0 & 0 & 0 & \hl{2} & 0 & 0 & 0 & 0 & 0 & 0 & 0 \\
\bottomrule
\end{tabular}
\\[4pt]
\begin{tabular}{>{\raggedright\arraybackslash}p{8.8em}>{\centering\arraybackslash}p{1.15em}>{\centering\arraybackslash}p{1.15em}>{\centering\arraybackslash}p{1.15em}>{\centering\arraybackslash}p{1.15em}>{\centering\arraybackslash}p{1.15em}>{\centering\arraybackslash}p{1.15em}>{\centering\arraybackslash}p{1.15em}>{\centering\arraybackslash}p{1.15em}>{\centering\arraybackslash}p{1.15em}>{\centering\arraybackslash}p{1.15em}>{\centering\arraybackslash}p{1.15em}>{\centering\arraybackslash}p{1.15em}>{\centering\arraybackslash}p{1.15em}>{\centering\arraybackslash}p{1.15em}>{\centering\arraybackslash}p{1.15em}>{\centering\arraybackslash}p{1.15em}}
\toprule
$n$ & \hl{17} & 18 & 19 & 20 & 21 & 22 & 23 & 24 & 25 & 26 & 27 & 28 & 29 & 30 & 31 & 32 \\ \midrule
$a(n)\bmod 4$ & \hl{2} & 0 & 0 & 0 & 0 & 0 & 0 & 0 & 0 & 0 & 0 & 0 & 0 & 0 & 0 & 0 \\
$\coeff{n}{(x{+}2xC)}\bmod 4$ & \hl{2} & 0 & 0 & 0 & 0 & 0 & 0 & 0 & 0 & 0 & 0 & 0 & 0 & 0 & 0 & 0 \\
\bottomrule
\end{tabular}
\\[4pt]
\begin{tabular}{>{\raggedright\arraybackslash}p{8.8em}>{\centering\arraybackslash}p{1.15em}>{\centering\arraybackslash}p{1.15em}>{\centering\arraybackslash}p{1.15em}>{\centering\arraybackslash}p{1.15em}>{\centering\arraybackslash}p{1.15em}>{\centering\arraybackslash}p{1.15em}>{\centering\arraybackslash}p{1.15em}>{\centering\arraybackslash}p{1.15em}>{\centering\arraybackslash}p{1.15em}>{\centering\arraybackslash}p{1.15em}>{\centering\arraybackslash}p{1.15em}>{\centering\arraybackslash}p{1.15em}>{\centering\arraybackslash}p{1.15em}>{\centering\arraybackslash}p{1.15em}>{\centering\arraybackslash}p{1.15em}>{\centering\arraybackslash}p{1.15em}}
\toprule
$n$ & \hl{33} & 34 & 35 & 36 & 37 & 38 & 39 & 40 & 41 & 42 & 43 & 44 & 45 & 46 & 47 & 48 \\ \midrule
$a(n)\bmod 4$ & \hl{2} & 0 & 0 & 0 & 0 & 0 & 0 & 0 & 0 & 0 & 0 & 0 & 0 & 0 & 0 & 0 \\
$\coeff{n}{(x{+}2xC)}\bmod 4$ & \hl{2} & 0 & 0 & 0 & 0 & 0 & 0 & 0 & 0 & 0 & 0 & 0 & 0 & 0 & 0 & 0 \\
\bottomrule
\end{tabular}
\end{center}

\subsection{What this paper proves}
The \emph{Catalan numbers} $1,1,2,5,14,42,\dots$ (OEIS A000108) arise throughout combinatorics;
for example, they count the full binary trees with $1,2,3,4,5,\dots$ leaves. Their generating function is
$C(x)=x+x^2+2x^3+5x^4+14x^5+\cdots$, the only series without constant term that satisfies $C=x+C^2$, and we write $c_n$
for its coefficient of $x^n$.

\textbf{Main result} (Theorem~\ref{thm:main}). \emph{Modulo $4$, the series $A(x)$ equals $x+2x\,C(x)$. In other
words, $a(n)-2c_{n-1}$ is divisible by $4$ for every $n>1$.}

The number $2c_{n-1}$ leaves remainder $2$ on division by $4$ when $c_{n-1}$ is odd, and remainder $0$ when $c_{n-1}$ is
even. A classical fact (Lemma~\ref{lem:catalan}) says that $c_j$ is odd exactly when $j$ is a power of $2$. Together
these prove Hanna's conjecture (Corollary~\ref{cor:conj}).

\subsection{What was known}
The OEIS entry records the defining equation, an equivalent equation for the compositional inverse of $A$, a formula
for $A$ as an infinite product, and numerical values of $A$ as a function. The parity of the Catalan numbers is classical;
see for example~\cite{deutschsagan}. We are not aware of an earlier proof of the conjecture.

\subsection{Motivation}
Mahler-type equations have a long history in number theory: Mahler introduced them to prove that certain numbers are
transcendental, and they are closely connected with automatic sequences, which are generated by finite
automata~\cite{allouche}. The equation here is not of the classical form, because of the extra term $2x^3$; but
modulo $2$ that term disappears, and the equation becomes the classical $A(x^2)=x\,A(x)$. The Catalan numbers, which the equation does not mention, describe its solution
modulo $4$. The proof finds a series that solves the equation modulo $4$ and then uses the uniqueness of the solution;
the same method is used in other papers of this collection.

\subsection{How to read this paper}
Section~\ref{sec:background} explains the notions we need: power series, substitution, remainders, the Catalan numbers
modulo $2$, and a lemma about doubling modulo $4$. Section~\ref{sec:results} states the results precisely, and
Section~\ref{sec:proofs} proves them. Section~\ref{sec:checks} describes how the results were checked, by computer and
by a formal proof in Lean. Section~\ref{sec:discussion} discusses the scope of the result and an open question.
\section{Background}\label{sec:background}
This section collects the definitions and facts used later. Readers who know formal power series and congruences can skip it and return when needed.
\subsection{Formal power series}
A \emph{formal power series} is an expression
\[
U(x)=u_0+u_1x+u_2x^2+u_3x^3+\cdots ,
\]
determined by its sequence of coefficients $u_0,u_1,u_2,\dots$; see~\cite{niven,stanley}. No number is ever
substituted for $x$, so questions of convergence do not arise. We write $\coeff{n}{U}=u_n$ for the coefficient of $x^n$ in $U$; two series are equal
when all their coefficients are equal.

Series are added coefficient by coefficient and multiplied like polynomials:
\[
\coeff{n}{(UV)}=u_0v_n+u_1v_{n-1}+\cdots+u_nv_0 .
\]
Each coefficient of a product is a finite sum. For example, $(1-x)(1+x+x^2+x^3+\cdots)=1$, so $1+x+x^2+\cdots$ is
the \emph{multiplicative inverse} of $1-x$, written $\frac{1}{1-x}$. More generally, every series $D$ with constant term
$1$ has a unique multiplicative inverse $\frac1D$, whose coefficients are determined one at a time by $D\cdot\frac1D=1$
using only addition, subtraction and multiplication. Whenever a quotient of power series appears in this paper, its
denominator has constant term $1$, so the quotient is again a power series.

A \emph{commutative ring} is a set with an addition, a subtraction and a multiplication that satisfy the usual rules
of arithmetic; examples are $\Z$, $\Q$ and the integers modulo $m$. For a commutative ring $R$ we write $\ps{R}$ for
the set of power series with coefficients in $R$; for instance, writing $\F_2$ for the integers modulo $2$, $\ps{\F_2}$
consists of the series with coefficients modulo $2$. We write
$U(0)$ for the constant term $u_0$ of $U$, and $U=O(x^d)$ when the coefficients of $1,x,\dots,x^{d-1}$ in $U$ are all
zero. Thus $U-V=O(x^d)$ means that $U$ and $V$ have the same coefficient of $x^n$ for every $n<d$. A series with
constant term $0$ and coefficient of $x$ equal to $1$, that is, a series of the form $x+u_2x^2+u_3x^3+\cdots$, is
called \emph{normalized}.

\subsection{Substituting one series into another}
If $V(x)$ has constant term $0$, we can substitute $V$ for $x$ in $U$:
\[
U(V(x))=u_0+u_1V(x)+u_2V(x)^2+u_3V(x)^3+\cdots .
\]
This is called \emph{composition}, also written $U\circ V$. It makes sense because $V(x)^k$ has no terms below $x^k$,
so only the finitely many terms with $k\le n$ affect the coefficient of $x^n$. For example, if $U=x+x^2$ and $V=x^2$, then
$U(V(x))=x^2+x^4$. Composition is associative: $U\circ(V\circ W)=(U\circ V)\circ W$.

Substituting a series into itself repeatedly gives its \emph{iterates}: $U^{\circ2}(x)=U(U(x))$,
$U^{\circ3}(x)=U(U(U(x)))$, and so on, with $U^{\circ0}(x)=x$ and $U^{\circ1}=U$. They must not be confused with
\emph{powers} $U^2=U\cdot U$, $U^3=U\cdot U\cdot U$. (The OEIS sometimes writes $A^k(x)$ for the $k$-th iterate. In
this paper the small circle always means an iterate, and a plain exponent always means a power.)

\subsection{Remainders and congruences}
For an integer $m\ge2$ we write $a\equiv b\pmod m$ (``$a$ is congruent to $b$ modulo $m$'') when $a-b$ is divisible by
$m$, that is, when $a$ and $b$ leave the same remainder after division by $m$. The remainder of $a$ on division by
$m$ is called its \emph{residue} modulo $m$; it is one of $0,1,\dots,m-1$.

Addition and multiplication are compatible with congruences, so the residues $0,1,\dots,m-1$, added and multiplied
modulo $m$, form a ring, written $\Z/m\Z$. When $m=p$ is a prime we also write $\F_p$. In $\F_2=\{0,1\}$ we have $1+1=0$,
so addition and subtraction coincide.

For power series with integer coefficients, $U\equiv V\pmod m$ means that \emph{every} coefficient of $U-V$ is
divisible by $m$. Equivalently, the two series become equal when every coefficient is replaced by its residue; we call
this \emph{reducing} the series modulo $m$. The key fact is the following.

\emph{Reduction modulo $m$ is compatible with all the operations above.} To reduce a sum, a product, a substitution, or
an inverse, we may reduce the ingredients first and then compute in $\Z/m\Z$. This holds because every
coefficient of the result is obtained from finitely many coefficients of the ingredients by adding, subtracting and
multiplying, and these operations are compatible with remainders.

\para{Rings.} Some statements below hold over every commutative ring, for example $\Z$, $\Q$ and $\Z/m\Z$. No
division is used in them; this matters because in $\Z/m\Z$ division is not always possible (for instance, $2$ has no inverse in
$\Z/4\Z$).

\subsection{The method}
The main proofs follow three steps.
\begin{enumerate}
\item \emph{Uniqueness.} Show that the defining equation determines the coefficients one at a time, using only
addition, subtraction and multiplication. Then the equation has exactly one solution of the required shape, in every
commutative ring, including $\Z/4\Z$.
\item \emph{Reduction.} Reducing the integer solution modulo $4$ gives a solution of the reduced equation,
because reduction is compatible with all operations.
\item \emph{A candidate solution.} Find an explicit series that satisfies the reduced equation. By uniqueness, it is
the reduction of the integer solution, and its coefficients can be read off.
\end{enumerate}
Finding the candidate in step~3 is usually the main difficulty.
\subsection{The Frobenius identity}
\begin{lemma}[Frobenius identity]\label{lem:frobenius}
Let $p$ be a prime. For every power series $U$ with coefficients in $\F_p$, we have $U(x)^p=U(x^p)$.
\end{lemma}

\begin{proof}
In $\F_p$, the binomial theorem simplifies to $(a+b)^p=a^p+b^p$, because every binomial coefficient
$\binom pi$ with $0<i<p$ is divisible by $p$. Also $c^p=c$ for every $c\in\F_p$ (Fermat's little theorem). So
$\bigl(\sum_nu_nx^n\bigr)^p=\sum_nu_n^px^{pn}=\sum_nu_nx^{pn}=U(x^p)$. The first equality holds for polynomials by
repeated use of $(a+b)^p=a^p+b^p$; it holds for power series because each coefficient of $U(x)^p$ involves only finitely
many coefficients of $U$.
\end{proof}

For $p=2$ this says that squaring a series modulo $2$ doubles every exponent: for instance,
$(x+x^3)^2\equiv x^2+x^6\pmod2$. In particular, the square of any series has, modulo $2$, no terms of odd degree.

\subsection{The Catalan numbers modulo 2}
\begin{lemma}\label{lem:catalan}
In every commutative ring the equation $U=x+U^2$ has exactly one solution with $U(0)=0$. Over the integers it is the
Catalan series $C(x)=x+x^2+2x^3+5x^4+14x^5+\cdots$, and
\[
C(x)\equiv x+x^2+x^4+x^8+\cdots=\sum_{k\ge0}x^{2^k}\pmod 2 .
\]
So $\coeff{n}{C}$ is odd exactly when $n$ is a power of $2$.
\end{lemma}

\begin{proof}
The coefficient of $x^n$ in $U^2$ involves only $u_1,\dots,u_{n-1}$, so the equation determines $u_n$ one at a time:
$u_1=1$ and $u_n=\sum_{i=1}^{n-1}u_iu_{n-i}$. Over $\F_2$, the series $S=\sum_kx^{2^k}$ satisfies
$x+S^2=x+\sum_kx^{2^{k+1}}=S$ by Lemma~\ref{lem:frobenius}, so by uniqueness $C\equiv S\pmod2$.
\end{proof}
The parity of the Catalan numbers is classical; see for example~\cite{deutschsagan}. (Here the Catalan number
$\frac1{n}\binom{2n-2}{n-1}$ appears as the coefficient of $x^n$, so ``$n$ a power of $2$'' corresponds to the familiar
statement that the $k$-th Catalan number is odd exactly when $k+1$ is a power of $2$.)

\subsection{Doubling modulo 4}
\begin{lemma}\label{lem:double}
Let $U,V$ be series with coefficients in $\Z/4\Z$. Then $2U=2V$ if and only if $U$ and $V$ agree modulo $2$.
Consequently, if $P$ and $Q$ have constant term $0$ and $P\equiv Q\pmod 2$, then $2H(P)=2H(Q)$ in $\Z/4\Z$ for every
series $H$.
\end{lemma}

\begin{proof}
For residues $u,v$ modulo $4$, $2u=2v$ means that $4$ divides $2(u-v)$, that is, $2$ divides $u-v$. Apply this to
every coefficient. For the second claim, $H(P)$ and $H(Q)$ agree modulo $2$ because reduction is compatible with
substitution.
\end{proof}
In other words, \emph{twice a series, modulo $4$, depends only on the series modulo $2$.} This lets us use information modulo
$2$ to get information modulo $4$, without ever dividing by $2$ in $\Z/4\Z$, which is impossible.

\section{Results}\label{sec:results}

Throughout, $C$ is the Catalan series, the unique normalized series with $C=x+C^2$ (Lemma~\ref{lem:catalan}), and
$c_n=\coeff{n}{C}$.

\begin{theorem}\label{thm:main}
\begin{enumerate}[label=(\alph*)]
\item Equation~\eqref{eq:def} has exactly one normalized solution $A$ with integer coefficients (and, more generally,
exactly one with coefficients in any commutative ring). The sequence A389536 is $a(n)=\coeff{n}{A}$.
\item Modulo $4$,
\[
A(x)\equiv x+2x\,C(x)\pmod 4 .
\]
\end{enumerate}
\end{theorem}

\begin{corollary}[Hanna's conjecture]\label{cor:conj}
For every $n>1$: if $n=2^k+1$ for some $k\ge0$, then $a(n)\equiv2\pmod4$; otherwise $a(n)$ is divisible by $4$.
\end{corollary}

\begin{proof}
For $n>1$ the coefficient of $x^n$ in $x+2xC(x)$ is $2c_{n-1}$, so Theorem~\ref{thm:main}(b) gives
$a(n)\equiv2c_{n-1}\pmod4$. This is $2$ modulo $4$ if $c_{n-1}$ is odd, and $0$ if it is even. By
Lemma~\ref{lem:catalan}, $c_{n-1}$ is odd exactly when $n-1$ is a power of $2$.
\end{proof}

For example, $c_8=429$ is odd, and indeed $a(9)=202=4\cdot50+2$; while $c_9=1430$ is even, and $a(10)=356=4\cdot89$.

\section{Proofs}\label{sec:proofs}

\para{The idea in brief.} Comparing coefficients in \eqref{eq:def} gives a formula for each coefficient in terms of
earlier ones, so there is exactly one solution with $a(1)=1$. For the congruence we take the candidate $D=x+2xC$ and check
that it satisfies \eqref{eq:def} modulo $4$. Modulo $2$ the inner series $x^2+2x^3$ is $x^2$, and the Catalan series
satisfies $C(x^2)\equiv C(x)^2\pmod 2$. Because $C$ appears in $D$ multiplied by $2$, information modulo $2$ is enough
(Lemma~\ref{lem:double}), and the identity $C^2=C-x$ finishes the calculation.

\begin{proof}[Proof of Theorem~\ref{thm:main}]
\emph{(a) Existence and uniqueness.} Write $P(x)=x^2+2x^3=x^2(1+2x)$, so that $P^j=x^{2j}(1+2x)^j$. By the binomial
theorem, the coefficient of $x^{n+1}$ in $P^j$ is $\binom{j}{n+1-2j}2^{\,n+1-2j}$ (and $0$ unless
$2j\le n+1\le3j$). So comparing the coefficients of $x^{n+1}$ on the two sides of \eqref{eq:def} gives
\[
a(n)=\sum_{j}\binom{j}{n+1-2j}\,2^{\,n+1-2j}\,a(j),
\]
where the sum runs over the $j$ with $2j\le n+1\le 3j$. For $n\ge2$ all these $j$ satisfy $j\le\frac{n+1}2<n$, so this
formula computes $a(n)$ from earlier coefficients. For $n=1$ it reads $a(1)=a(1)$, which is no condition, and for $n=0$
it gives $a(0)=0$. Hence, once $a(1)=1$ is fixed, there is exactly one solution, and it is computed using only additions
and multiplications. This works in every commutative ring, and reducing the integer solution modulo $4$ gives the solution
over $\Z/4\Z$. For example, $a(2)=\binom11\cdot2\cdot a(1)=2$, $a(3)=\binom20\cdot a(2)=2$, and
$a(4)=\binom21\cdot2\cdot a(2)=8$.

\emph{(b) Modulo $4$.} By part (a) it suffices to check that the normalized series $D=x+2x\,C$ satisfies
\eqref{eq:def} over $\Z/4\Z$. Substituting $P=x^2+2x^3$ into $D$ gives
\[
D(P)=P+2P\,C(P)=x^2+2x^3+2x^2\,C(P)+4x^3\,C(P)=x^2+2x^3+2x^2\,C(P)\quad\text{in }\Z/4\Z .
\]
To simplify $2\,C(P)$, note that $P\equiv x^2\pmod2$. By Lemma~\ref{lem:double}, $2\,C(P)=2\,C(x^2)$ in $\Z/4\Z$. By
Lemma~\ref{lem:frobenius}, $C(x^2)\equiv C(x)^2\pmod 2$, so again by Lemma~\ref{lem:double}, $2\,C(x^2)=2\,C^2$. Finally
$C^2=C-x$. Therefore $2x^2\,C(P)=2x^2(C-x)=2x^2C-2x^3$, and
\[
D(P)=x^2+2x^3+2x^2C-2x^3=x^2+2x^2C=x\,D .
\]
So $D$ satisfies \eqref{eq:def} over $\Z/4\Z$, and by uniqueness $A\equiv D=x+2xC\pmod4$.
\end{proof}

\section{Checking the results}\label{sec:checks}

A mathematical proof can contain mistakes, so the results were checked in two independent ways.

\para{Numerical check.}
We solved \eqref{eq:def} coefficient by coefficient with exact integer arithmetic. This reproduces all 34 terms
printed in the OEIS entry, and the rule of Corollary~\ref{cor:conj} holds for every $n$ with $1<n<130$.
The computation was done with the PARI/GP computer algebra system, using code that was written separately from the
formal proofs (file \texttt{verification/independent\_checks.gp} in the accompanying repository~\cite{repo}).
A computation can only test finitely many cases, so it is a check on the definitions and on the indexing, not a proof.

\para{Formal proof in Lean.}
Lean~\cite{lean} is a \emph{proof assistant}: a program in which mathematical statements and their proofs are written
in a precise formal language, and every step is checked mechanically down to the basic rules of logic. The theorems of Section~\ref{sec:results} are proved in Lean~4 (version 4.35.0) on top of
Mathlib~\cite{mathlib}, Lean's library of formalized mathematics. The statements are collected in the file \texttt{lean/Entries/A389536.lean} of the repository,
in the namespace \texttt{OEIS.A389536}. There the Catalan series appears as \texttt{scaledSolution 0}, the solution of a more general equation which, for the parameter value $0$, says $C(x-x^2)=x$; the series with this property is the same as the solution of $C=x+C^2$. The proofs are in the module \texttt{lean/Families/MahlerParity.lean}.
The table lists where each result of this paper is proved.
\begin{center}\small\renewcommand{\arraystretch}{1.25}
\begin{tabular}{>{\raggedright\arraybackslash}p{0.52\linewidth}>{\raggedright\arraybackslash}p{0.4\linewidth}}
\toprule
Statement in this paper & Lean theorem(s)\\
\midrule
Theorem~\ref{thm:main}(a) for integer coefficients; $A$ is normalized and satisfies \eqref{eq:def} & \texttt{mahler\_\allowbreak{}existence}, \texttt{mahler\_\allowbreak{}definition}, \texttt{mahler\_\allowbreak{}source\_\allowbreak{}unique} \\
Theorem~\ref{thm:main}(b): $A\equiv x+2xC\pmod4$ & \texttt{mahler\_\allowbreak{}mod\_\allowbreak{}four} \\
Corollary~\ref{cor:conj} & \texttt{mahler\_\allowbreak{}residue}, \texttt{mahler\_\allowbreak{}four\_\allowbreak{}dvd\_\allowbreak{}iff} \\
\bottomrule
\end{tabular}
\end{center}
The Lean project was rebuilt from source, and the axioms on which each of these theorems depends were listed with
Lean's \texttt{\#print axioms} command. Every one depends only on the three standard axioms of Lean's logic
(\texttt{propext}, \texttt{Classical.choice}, \texttt{Quot.sound}); no step is left unproved. A formal proof certifies
exactly the statement that Lean checks. Whether that statement says what the OEIS entry means is a separate question,
so the Lean definitions were also compared with the definition in the OEIS entry
(Section~\ref{sec:sequence}). The principal theorems are stated once more, using only Mathlib's definitions, in the
file \texttt{lean/Challenges/Catalan.lean}, and Lean's Comparator tool confirms that the
proved theorems have exactly these statements.

\section{Discussion}\label{sec:discussion}

\para{Scope.} Theorem~\ref{thm:main} determines the remainder of every coefficient $a(n)$ modulo $4$, so it settles the
conjecture of the entry completely. The statements in the entry about the radius of convergence and about the values of
$A$ as a function are of a different nature (they concern the size of the coefficients, not their remainders) and are not
addressed here.

\para{Modulo 8.} A computer calculation suggests a more intricate rule modulo $8$. For every $n$ with $1<n<129$ we found:
$a(n)\equiv2\pmod8$ when $n-1$ is a power of $2$; $a(n)\equiv4\pmod8$ when $n-1=2^i+2^j$ with $i\ge j+2$ (that is,
$n=6,10,11,18,19,21,34,\dots$); and $a(n)\equiv0\pmod8$ otherwise. We have not tried to prove this.

\section*{About this paper}\label{sec:about}
\para{How the work was done.}
This paper is one of 29 papers on conjectures from the OEIS that came out of an AI research effort directed by
Dev~D.~Goyal~\cite{repo}. The roles were as follows.
\begin{itemize}
\item \emph{Dev~D.~Goyal} started the project, chose which problems to pursue, asked for proofs, formal verification and
revisions, and decided what to publish. Dev~D.~Goyal did not supply proof strategies or intermediate arguments.
\item \emph{ChatGPT} (OpenAI), using the GPT-6 model with Pro reasoning, worked autonomously from those prompts. It
searched the OEIS for suitable conjectures, found the proofs, wrote the Lean~4 formalization, and wrote a first draft
of this paper.
\item \emph{Claude Opus~5.5}, used with Claude Code, checked the work independently. It rebuilt the Lean project from source and confirmed that the theorems listed in Section~\ref{sec:checks} rely only on
Lean's standard axioms; compared the Lean definitions with the OEIS definition; recomputed the sequence with separately
written PARI/GP code, compared it with the terms published in the OEIS, and tested every claim of the paper on many
terms; and read every step of the proofs. It then wrote the present version of the paper, expanding the
explanations and, in places, writing out an argument in more detail or by a more elementary route. The theorems of
Section~\ref{sec:results} and the proof strategies are ChatGPT's. The Discussion section, including the short arguments
in it, was written by Claude Opus~5.5; its observations were checked by computer and are not part of the Lean
formalization.
\end{itemize}
Dev~D.~Goyal has reviewed the proofs and is responsible for this paper; this review is not an independent expert
validation.

\para{Review status.} This paper has not been peer reviewed, and it has not been submitted to a journal or to the OEIS.

\para{Priority.} An OEIS entry that labels a statement as a conjecture shows only that its contributor did not know
a proof; it does not show that no proof existed. During the research only limited searches for earlier work were
made. No claim of priority is made.

\para{Corrections.} Please report any error in the issue tracker of the repository~\cite{repo}, where the Lean sources and the numerical checks are also
available.

\section*{Acknowledgements}
The sequence, its formulas and the conjecture are due to Paul~D.~Hanna~\cite{oeis}.

\begin{thebibliography}{99}\small
\raggedright
\bibitem{oeis} P.~D.~Hanna, entry A389536 (G.f.\ $A(x)$ satisfies $A(x)=A(x^2+2x^3)/x$), \emph{The On-Line Encyclopedia of Integer Sequences}, October 2025. \url{https://oeis.org/A389536}. Text as of October 2026.
\bibitem{deutschsagan} E.~Deutsch and B.~E.~Sagan, Congruences for Catalan and Motzkin numbers and related sequences, \emph{Journal of Number Theory} 117 (2006), 191--215.
\bibitem{allouche} J.-P.~Allouche and J.~Shallit, \emph{Automatic Sequences}, Cambridge University Press, 2003.
\bibitem{oeisdb} N.~J.~A.~Sloane and The OEIS Foundation Inc., \emph{The On-Line Encyclopedia of Integer Sequences}. \url{https://oeis.org}.
\bibitem{stanley} R.~P.~Stanley, \emph{Enumerative Combinatorics}, Volume~2, Cambridge University Press, 1999. (Chapter~5 treats formal power series, composition and compositional inverses.)
\bibitem{niven} I.~Niven, Formal power series, \emph{American Mathematical Monthly} 76 (1969), 871--889.
\bibitem{lean} L.~de~Moura and S.~Ullrich, The Lean~4 theorem prover and programming language, in \emph{Automated Deduction -- CADE 28}, Lecture Notes in Computer Science 12699, Springer, 2021, 625--635.
\bibitem{mathlib} The mathlib Community, The Lean mathematical library, in \emph{Proceedings of the 9th ACM SIGPLAN International Conference on Certified Programs and Proofs (CPP 2020)}, 367--381. Version used: commit \texttt{56a005855ea1}.
\bibitem{repo} D.~D.~Goyal, \emph{Proofs and corrections of OEIS conjectures}, repository, 2026. \url{https://github.com/dvgyl/oeis-proofs}.
\end{thebibliography}
\end{document}
